% Standalone preview. Copy only the subsection into the system model.
% Citation keys match docs/references.bib.
\documentclass[10pt,twocolumn]{article}
\usepackage[a4paper,margin=19mm]{geometry}
\usepackage{amsmath,amssymb}
\usepackage[hidelinks]{hyperref}
\begin{document}

\subsection{Power Model and Energy Model}
\label{subsec:energy_model}

DVFS adjusts the execution frequency of task $i$ through the ratio
$\rho_i$. For its assigned server $m_i$, the execution frequency is
\begin{equation}
 f_i=\rho_iF_{m_i},\qquad \underline\rho_{m_i}\leq\rho_i\leq1,
 \label{eq:task_frequency}
\end{equation}
where $F_m$ and $\underline\rho_m$ are the maximum frequency and minimum
supported frequency ratio of server $m$, respectively; the minimum
frequency is $F_m^{\min}=\underline\rho_m F_m$.
Following the frequency-dependent power models
in~\cite{xie2017energy,zhu2006energy}, we adopt a cubic active-power model:
\begin{equation}
 P_i=P_{m_i}^{\mathrm{ind}}+P_{m_i}^{\mathrm{dyn}}\rho_i^3,
 \label{eq:task_power}
\end{equation}
where $P_m^{\mathrm{ind}}$ is frequency-independent active power and
$P_m^{\mathrm{dyn}}$ is dynamic power at the maximum frequency.
The frequency remains constant during task execution. With execution
time $ET_i=L_i/f_i$, the task's computation energy is
\begin{equation}
 E_i=P_iET_i
 =\frac{L_i}{F_{m_i}}\left(
 \frac{P_{m_i}^{\mathrm{ind}}}{\rho_i}
 +P_{m_i}^{\mathrm{dyn}}\rho_i^2\right).
 \label{eq:task_energy}
\end{equation}

Reducing frequency lowers dynamic power but prolongs execution and
therefore increases the frequency-independent energy component.
Minimizing $E_i$ with respect to $\rho_i$ over the supported interval gives
the energy-efficient operating ratio of server $m$:
\begin{equation}
 \rho_m^{\mathrm{ee}}=\min\!\left\{1,\max\!\left\{
 \underline\rho_m,
 \left(\frac{P_m^{\mathrm{ind}}}{2P_m^{\mathrm{dyn}}}\right)^{1/3}
 \right\}\right\}.
 \label{eq:rhoee}
\end{equation}
This ratio depends on server power characteristics, not task workload;
meeting timing and reliability requirements may require a higher
execution frequency.
The total computation energy of workflow $w$ is
\begin{equation}
 E_w=\sum_{i\in\mathcal N}E_i.
 \label{eq:workflow_energy}
\end{equation}
Only active computation energy is considered; idle and communication
energy are excluded.

% Preview-only bibliography. Use references.bib in the full manuscript.
\begin{thebibliography}{2}
\bibitem{xie2017energy}
G. Xie, G. Zeng, X. Xiao, R. Li, and K. Li,
``Energy-efficient scheduling algorithms for real-time parallel
applications on heterogeneous distributed embedded systems,''
\emph{IEEE Trans. Parallel Distrib. Syst.}, vol. 28, no. 12,
pp. 3426--3442, 2017, doi: 10.1109/TPDS.2017.2730876.
\bibitem{zhu2006energy}
D. Zhu and H. Aydin, ``Energy management for real-time embedded systems
with reliability requirements,'' in \emph{Proc. IEEE/ACM ICCAD},
2006, pp. 528--534, doi: 10.1109/ICCAD.2006.320169.
\end{thebibliography}
\end{document}
